\section{Conclusion}\label{Sec:Conc}

We reassess BKK's findings on international business cycles using 23 additional years of data. Standard deviations and within-country comovements change relatively little, while cross-country correlations decline for most variables.

Using the United States as a benchmark produces results close to those obtained from all country pairs. In this sample, its distinctive size and degree of openness do not substantially distort the summary of international comovements.

Our estimates lie between BKK's and ACZ's. The lower correlations reported by ACZ may owe more to their larger country sample than to the length of their time series.

The quantity anomaly persists. Output correlation exceeds consumption correlation in 53 of the 55 country pairs. All output correlations are positive, in contrast to the negative correlation predicted by the benchmark model. No consumption correlation reaches the model's prediction of 0.88.

The overall decline in cross-country correlations does not resolve the anomaly because it concerns the relative magnitudes of output and consumption correlations. Both decline in our sample, leaving their ranking largely unchanged.